\documentclass[11pt]{article}

\usepackage[margin=0.9in]{geometry}
\usepackage{amsmath,amssymb,amsthm,mathtools,booktabs,enumitem,microtype}
\usepackage{xcolor}
\usepackage[colorlinks=true,linkcolor=blue!50!black,urlcolor=blue!50!black,citecolor=blue!50!black]{hyperref}

\theoremstyle{plain}
\newtheorem{theorem}{Theorem}
\newtheorem{lemma}[theorem]{Lemma}
\newtheorem{proposition}[theorem]{Proposition}
\newtheorem{corollary}[theorem]{Corollary}

\theoremstyle{definition}
\newtheorem{definition}[theorem]{Definition}

\theoremstyle{remark}
\newtheorem{remark}{Remark}

\DeclareMathOperator{\rank}{rank}
\DeclareMathOperator{\Sym}{Sym}
\DeclareMathOperator{\diag}{diag}
\DeclareMathOperator{\spn}{span}

\title{\bfseries Refutation of Global Identifiability for Deep Polynomial Neural Networks}
\author{}
\date{\vspace{-1.6em}%
\small
\begin{tabular}{@{}c@{}}
Target paper: K.~Usevich, R.~Borsoi, C.~D\'erand, and M.~Clausel,\\
\emph{Identifiability of Deep Polynomial Neural Networks},\\
\href{https://arxiv.org/abs/2506.17093}{arXiv:2506.17093} (v3, NeurIPS 2025), Section~4.3, Conjecture~\texttt{conj:global\_identifiability}\\[0.35em]
November 2026
\end{tabular}}

\begin{document}

\maketitle

\begin{abstract}
We refute Conjecture~\texttt{conj:global\_identifiability} of Usevich, Borsoi, D\'erand, and Clausel (\href{https://arxiv.org/abs/2506.17093}{arXiv:2506.17093} v3, NeurIPS~2025) over both $\mathbb{R}$ and $\mathbb{C}$.
A generic two-layer homogeneous polynomial neural network (hPNN) with width tuple $(4, 7, 7)$ and activation degree $r_1 = 2$ admits \emph{exactly $8$} inequivalent representations modulo hidden-neuron permutation and scaling over both $\mathbb{C}$ and $\mathbb{R}$, arising from the classical Cayley octad of three quadrics in $\mathbb{P}^3$.
Lifting this shallow fiber to every depth $L \ge 3$ in the family $(4, 7, \dots, 7)$ with all activation degrees $r_j = 2$, every reduced two-layer block in the paper's Localization Theorem equals $((4, 7, 7), 2)$ and is therefore finitely identifiable; hence the deep generic fiber is finite and contains \emph{at least $8$} inequivalent representations over both fields.
We also verify global identifiability of the width-two formats $(2,2,2)$ and $(2,2,2,2)$ at degree~$2$ and note a minor typo in Appendix~E of arXiv:2506.17093~v3.
\end{abstract}

\section{Setup and Reduction of the Two-Layer Fiber}

An $L$-layer bias-free homogeneous polynomial neural network (hPNN) of architecture $(\mathbf{d}, \mathbf{r}) = ((d_0, \dots, d_L), (r_1, \dots, r_{L-1}))$ over $\mathbb{K} \in \{\mathbb{R}, \mathbb{C}\}$ defines the polynomial map
\[
F_W = W_L \circ \rho_{r_{L-1}} \circ W_{L-1} \circ \cdots \circ \rho_{r_1} \circ W_1,
\]
where $W_j \in \mathbb{K}^{d_j \times d_{j-1}}$ and $\rho_r$ acts coordinatewise by $z \mapsto z^r$.
Two weight tuples are equivalent if they differ by hidden-neuron permutations and coordinatewise rescalings multiplying row $i$ of $W_j$ by $s \neq 0$ and column $i$ of $W_{j+1}$ by $s^{-r_j}$.
An architecture is \emph{finitely identifiable} (resp.\ \emph{globally identifiable}) if a generic weight tuple (off a proper Zariski-closed set, hence Lebesgue-almost everywhere over $\mathbb{R}$) has finitely many (resp.\ a unique) equivalence classes in its fiber.

The paper's Localization Theorem (\texttt{thm:localization\_theorem}) proves that if every reduced two-layer architecture $((\tilde{d}_{\ell-1}, d_\ell, d_{\ell+1}), r_\ell)$ with $\tilde{d}_{\ell-1} := \min_{0 \le i < \ell} d_i$ is finitely identifiable ($\ell = 1, \dots, L-1$), then the $L$-layer architecture $(\mathbf{d}, \mathbf{r})$ is finitely identifiable.
Conjecture~\texttt{conj:global\_identifiability} asserts that under the exact same hypotheses, the $L$-layer architecture is \emph{globally} identifiable.
(The paper's separate Kruskal criterion, Proposition~\texttt{cor:gen\_generic}, requires $r \ge 7/3$ on $(4,7,7)$, so finite identifiability at $r = 2$ must be proved directly.)

For $((4, 7, 7), 2)$, write $F(x) = B(\ell_1(x)^2, \dots, \ell_7(x)^2)^\top$, where $\ell_i(x) = a_i \cdot x$ ($a_i \in \mathbb{K}^4$) and $B \in \mathbb{K}^{7 \times 7}$. Whenever $\det B \neq 0$ and $\ell_1^2, \dots, \ell_7^2$ are linearly independent, the output quadric span $S := \spn\{F_1, \dots, F_7\} \subset \Sym^2(\mathbb{K}^4)$ has dimension $7$.
Any representation $F(x) = \tilde{B}(\tilde{\ell}_1(x)^2, \dots, \tilde{\ell}_7(x)^2)^\top$ (even allowing degenerate weights a priori) satisfies $\spn\{\tilde{\ell}_1^2, \dots, \tilde{\ell}_7^2\} \supseteq S$ and has at most $7$ generators, so $\{\tilde{\ell}_1^2, \dots, \tilde{\ell}_7^2\}$ must be a basis of $S$; once the ordered linear forms are fixed, $\tilde{B}$ is uniquely determined.
Thus the fiber modulo equivalence is in bijection with \emph{unordered $7$-element sets of projective points $[a] \in \mathbb{P}^3$ whose squares $\ell_a^2$ form a basis of $S$}.

Identify $\ell_a^2$ in the $10$-element monomial basis $\{x_i^2, 2x_i x_j : 0 \le i < j \le 3\}$ with its quadratic Veronese vector $v(a) := (a_i a_j)_{0 \le i \le j \le 3} \in \mathbb{K}^{10}$. Each linear functional on $\mathbb{K}^{10}$ evaluates on $v(a)$ as a homogeneous quadric $Q(a)$, so
\begin{equation}\label{eq:annihilator}
\ell_a^2 \in S
\quad\Longleftrightarrow\quad
Q(a) = 0 \text{ for every quadric } Q \text{ vanishing on } [a_1], \dots, [a_7].
\end{equation}

\section{Exact Two-Layer Fiber Count via Cayley Octads}

The base locus of the net of quadrics through seven general points in $\mathbb{P}^3$ is a classical \emph{Cayley octad} of eight points~\cite[\S 7]{plaumann2011quartic}. We give a self-contained proof of the rank, transversality, genericity, and reality properties establishing the exact count~$8$.

\begin{theorem}[Exact fiber size $8$ for $((4, 7, 7), 2)$]\label{thm:shallow}
There exists a nonempty Zariski-open subset $\mathcal{U}$ of the weight space for $((4, 7, 7), 2)$ (whose complement has real Lebesgue measure zero) on which every network has \emph{exactly $8$} inequivalent representations over $\mathbb{C}$. Moreover, for every real weight tuple in $\mathcal{U}$, all $8$ equivalence classes are realized by real weight tuples.
\end{theorem}

\begin{proof}
\emph{Step~1 (Cube witness).} Consider the $8$ cube vertices $p_\epsilon := [1 : \epsilon_1 : \epsilon_2 : \epsilon_3] \in \mathbb{P}^3$, $\epsilon = (\epsilon_1, \epsilon_2, \epsilon_3) \in \{\pm 1\}^3$. On these $8$ points, the $10$ columns of the $8 \times 10$ Veronese evaluation matrix $V$ are the constant vector $1$ (on the four squares $a_k^2$) and the six sign characters $\epsilon_1, \epsilon_2, \epsilon_3, \epsilon_1\epsilon_2, \epsilon_1\epsilon_3, \epsilon_2\epsilon_3$. These $7$ characters are mutually orthogonal in $\mathbb{R}^8$, so $\rank(V) = 7$ and its left kernel is spanned by the eighth character $\epsilon_1\epsilon_2\epsilon_3$, giving the exact identity
\begin{equation}\label{eq:cube-syzygy}
\sum_{\epsilon \in \{\pm 1\}^3} \epsilon_1\epsilon_2\epsilon_3\,(x_0 + \epsilon_1 x_1 + \epsilon_2 x_2 + \epsilon_3 x_3)^2 = 0.
\end{equation}
Because every coefficient $\epsilon_1\epsilon_2\epsilon_3 \in \{\pm 1\}$ is nonzero, omitting \emph{any} single vertex leaves $7$ linearly independent Veronese rows. Hence the space of quadrics vanishing on any $7$ cube vertices has dimension $10 - 7 = 3$, spanned by $Q_j(a) = a_j^2 - a_0^2$ ($j = 1, 2, 3$). Their common zeros in $\mathbb{P}^3$ are precisely the $8$ cube vertices, all in the chart $a_0 = 1$, where the Jacobian $\partial(Q_1, Q_2, Q_3)/\partial(a_1, a_2, a_3) = \diag(2\epsilon_1, 2\epsilon_2, 2\epsilon_3)$ has rank~$3$ at all $8$ points.

\emph{Step~2 (Genericity over $\mathbb{C}$).} On the irreducible parameter space $(\mathbb{P}^3)^7$, the locus $U_0$ where the $7 \times 10$ Veronese matrix $V(a_1, \dots, a_7)$ has rank~$7$ is Zariski-open and contains the cube witness. On each $7 \times 7$ minor chart containing the witness, linear elimination yields a basis $(Q_1, Q_2, Q_3)$ of annihilating quadrics with regular coefficients. In the incidence variety
\[
\mathcal{Z} := \bigl\{((a_1, \dots, a_7), [a]) \in U_0 \times \mathbb{P}^3 : Q_1(a) = Q_2(a) = Q_3(a) = 0\bigr\},
\]
the subset where the $3 \times 4$ Jacobian of $(Q_1, Q_2, Q_3)$ in $a$ has rank $< 3$ is closed, so by projectivity of $\mathbb{P}^3$ its image in $U_0$ is Zariski-closed. Deleting it leaves a nonempty Zariski-open subset $U_1 \subseteq U_0$ containing the witness, over which every common zero of $(Q_1, Q_2, Q_3)$ is transverse and isolated. By projective B\'ezout's theorem, every fiber of $\mathcal{Z}|_{U_1} \to U_1$ consists of exactly $2^3 = 8$ distinct points.

Since $\mathcal{Z}|_{U_1} \to U_1$ is projective with finite fibers and full relative Jacobian rank, it is a finite \'etale cover of degree~$8$. The $7$ tautological sections $s_i(a_1, \dots, a_7) := [a_i]$ have disjoint open-and-closed images in $\mathcal{Z}|_{U_1}$. Their complement is therefore a degree-$1$ finite \'etale cover of $U_1$, giving an algebraic section $s_8 : U_1 \to \mathbb{P}^3$ for the residual eighth point $[a_8]$. For each $k \in \{1, \dots, 8\}$, the condition that the $7$ Veronese rows $\{v(s_j)\}_{j \neq k}$ have rank~$7$ is open on $U_1$ and holds at the cube witness; intersecting these $8$ open sets gives a nonempty Zariski-open set $U_2 \subseteq U_1$. Pulling $U_2$ back to nonzero rows of $W_1$ and intersecting with $\{\det B \neq 0\}$ defines the nonempty Zariski-open subset $\mathcal{U}$ of weight space.

\emph{Step~3 (Counting over $\mathbb{C}$ and $\mathbb{R}$).} Over $\mathcal{U}$,~\eqref{eq:annihilator} implies that $[a_1], \dots, [a_8]$ are the only projective square directions in $S$, and every $7$-element subset spans $S$. Distinct omissions yield distinct unordered $7$-sets of projective first-layer rows (an invariant of hidden-neuron permutations and rescalings), while fixing the $7$-set determines the network up to those symmetries. Hence the complex fiber has size exactly $\binom{8}{7} = 8$.

If $(W_1, B) \in \mathcal{U}$ is real, complex conjugation preserves the real quadrics $Q_1, Q_2, Q_3$ and hence permutes $\{[a_1], \dots, [a_8]\}$; fixing $[a_1], \dots, [a_7]$, it must fix $[a_8]$ as well. Thus $[a_8]$ has a real representative $a_8 \in \mathbb{R}^4 \setminus \{0\}$, and the unique change-of-basis matrices between the $8$ real square bases of $S$ are real. Since the complement of $\mathcal{U}$ is a proper algebraic subset, it has real Lebesgue measure zero.
\end{proof}

\section{Lift to Arbitrary Depth \texorpdfstring{$L \ge 3$}{L >= 3}}

\begin{theorem}[Counterexample at every depth $L \ge 3$]\label{thm:deep}
For every $L \ge 3$, consider the $L$-layer hPNN with widths $\mathbf{d} = (4, 7, \dots, 7) \in \mathbb{Z}^{L+1}$ and degrees $\mathbf{r} = (2, \dots, 2) \in \mathbb{Z}^{L-1}$.
\begin{enumerate}[label=\textup{(\arabic*)}]
\item Every reduced two-layer block $((\tilde{d}_{\ell-1}, d_\ell, d_{\ell+1}), r_\ell)$ ($\ell = 1, \dots, L-1$) equals $((4, 7, 7), 2)$ and is finitely identifiable by Theorem~\ref{thm:shallow}. All hypotheses of the Localization Theorem hold (see \texttt{thm:localization\_theorem}), so the $L$-layer architecture has a \emph{finite} generic fiber.
\item Over both $\mathbb{R}$ and $\mathbb{C}$, every generic weight tuple $(W_1, \dots, W_L)$ in this architecture admits \emph{at least $8$} inequivalent representations. Thus the architecture is not globally identifiable.
\end{enumerate}
\end{theorem}

\begin{proof}
Since $d_0 = 4$ and $d_1 = \cdots = d_L = 7$, we have $\tilde{d}_{\ell-1} = \min_{0 \le i < \ell} d_i = 4$ for all $\ell = 1, \dots, L-1$, so every reduced block is $((4, 7, 7), 2)$. Theorem~\ref{thm:shallow} supplies finite identifiability of each block, and the Localization Theorem implies that the $L$-layer generic fiber is finite.

For any weight tuple $(W_1, W_2, \dots, W_L)$ with $(W_1, W_2)$ in the Zariski-open set $\mathcal{U}$ of Theorem~\ref{thm:shallow}, let $(W_1^{(k)}, W_2^{(k)})$ ($k = 1, \dots, 8$, all real if $W_1, W_2$ are real) be the $8$ inequivalent shallow representations satisfying $W_2^{(k)} \rho_2(W_1^{(k)} x) = W_2 \rho_2(W_1 x)$ identically in $x$. Composing with $W_L \circ \rho_2 \circ \cdots \circ W_3 \circ \rho_2$ yields $8$ deep representations $(W_1^{(k)}, W_2^{(k)}, W_3, \dots, W_L)$ of the same map $F_W$. Because equivalence at any depth preserves the unordered set of projective rows of $W_1$, these $8$ tuples lie in pairwise distinct equivalence classes.
\end{proof}

\begin{remark}[Scope and relation to shallow identifiability theorems]
Theorem~\ref{thm:shallow} proves the \emph{exact} generic fiber count $8$ for $L = 2$, while Theorem~\ref{thm:deep} proves a \emph{finite} generic fiber of size \emph{at least $8$} for $L \ge 3$.
A strengthened conjecture requiring \emph{global} identifiability of every reduced two-layer block is not refuted by this construction.
In the shallow additive-decomposition literature, results upgrading finite identifiability (non-defectivity) to generic global identifiability---such as Casarotti--Mella~\cite{casarotti2022non} and Massarenti--Mella~\cite{massarenti2025bronowski}---explicitly isolate exceptional non-identifiable cases, of which the Cayley octad in $((4,7,7), 2)$ is an archetypal instance.
\end{remark}

\section{Width-Two Control Checks and Appendix E Typo}

We also verify the width-two quadratic architectures associated with the specific non-uniqueness example (\texttt{ex:no\_localization\_uniqueness}) in \href{https://arxiv.org/abs/2506.17093}{arXiv:2506.17093}~v3.

\begin{proposition}[Global identifiability of $((2,2,2), 2)$ and $((2,2,2,2), (2,2))$]
Both $((2, 2, 2), 2)$ and $((2, 2, 2, 2), (2, 2))$ are globally identifiable over $\mathbb{R}$ and $\mathbb{C}$.
\end{proposition}

\begin{proof}
For $((2,2,2), 2)$ with independent linear forms $u, v$, any square $(\alpha u + \beta v)^2 = \alpha^2 u^2 + 2\alpha\beta uv + \beta^2 v^2 \in \spn\{u^2, v^2\}$ satisfies $\alpha\beta = 0$, leaving only $u^2$ and $v^2$ up to scale.

For $((2,2,2,2), (2,2))$ with invertible weight matrices, let $q_1 = a u^2 + b v^2$ and $q_2 = c u^2 + d v^2$, so the output quartics span $\langle q_1^2, q_2^2 \rangle$. Any square $(Au^2 + Buv + Cv^2)^2$ in this span has vanishing odd-degree $u^3v$ and $uv^3$ coefficients, forcing $AB = BC = 0$. If $B = 0$, viewing $u^2, v^2$ as independent variables reduces to the shallow argument and forces $(Au^2 + Cv^2)^2$ to be proportional to $q_1^2$ or $q_2^2$. The sole remaining candidate is $(uv)^2$, which belongs to $\langle q_1^2, q_2^2 \rangle$ if and only if $a^2 d^2 - b^2 c^2 = 0$ (by eliminating $u^4$ and $v^4$). Generically $a^2 d^2 - b^2 c^2 \neq 0$, so $q_1, q_2$ and then $u, v$ are uniquely recovered modulo equivalence.
\end{proof}

\begin{remark}[Typo in Appendix~E of \href{https://arxiv.org/abs/2506.17093}{arXiv:2506.17093}~v3]
Example~\texttt{ex:no\_localization\_uniqueness} uses $(a, b, c, d) = (1, 1, 1, -1)$, lying on the exceptional locus $a^2 d^2 - b^2 c^2 = 0$ where $(uv)^2$ enters the quartic pencil. In Appendix~E (\texttt{app-new-E-proof-update.tex}), the sentence introducing this example prints $\mathbf{d} = (2, 2, 2), \mathbf{r} = (2, 2)$; the network in the example has three linear layers $(I_2, W_2, I_2)$ and two quadratic activations, so the intended format is $\mathbf{d} = (2, 2, 2, 2), \mathbf{r} = (2, 2)$.
\end{remark}

\section{Exact Computational Verification}

The script \texttt{solutions/deep\_poly\_ident/explore.py} performs exact rational verification in SymPy: starting from the output coefficient matrix, it computes the annihilating quadrics and solves them by lexicographic Gr\"obner elimination across disjoint affine charts of $\mathbb{P}^3$. On the cube witness and four seeded integer $(4,7,7)$ instances, it finds all $8$ rational projective directions, reconstructs all $\binom{8}{7} = 8$ bases per instance, and verifies all $40$ reconstructed coefficient matrices against their targets. On the cube instance it also independently expands the full quartic output of $(4,7,7,7)$ across all $8$ deep representations and verifies exact polynomial equality, alongside the $(2,2,2)$ and $(2,2,2,2)$ binary controls.

\begin{thebibliography}{9}
\bibitem{casarotti2022non}
A.~Casarotti and M.~Mella.
From non-defectivity to identifiability.
\emph{Journal of the European Mathematical Society}, 25(3):931--945, 2022.

\bibitem{massarenti2025bronowski}
A.~Massarenti and M.~Mella.
Bronowski's conjecture and the identifiability of projective varieties.
\href{https://arxiv.org/abs/2511.19703}{arXiv:2511.19703}, 2025.

\bibitem{plaumann2011quartic}
D.~Plaumann, B.~Sturmfels, and C.~Vinzant.
Quartic curves and their bitangents.
\emph{Journal of Symbolic Computation}, 46(6):712--733, 2011.

\bibitem{usevich2025identifiability}
K.~Usevich, R.~Borsoi, C.~D\'erand, and M.~Clausel.
Identifiability of deep polynomial neural networks.
\href{https://arxiv.org/abs/2506.17093}{arXiv:2506.17093} (v3, NeurIPS 2025), 2025.
\end{thebibliography}

\end{document}
